\section{Greek Sources: Nature, Convention, and Rational Judgment}

These sources do not contain the later Catholic taxonomy. They address antecedent questions: whether civic command is the final measure of justice; whether nature names rational order, biological regularity, or superior power; whether unwritten norms can judge local enactments; and how general rules should govern exceptional cases. Their vocabulary and genres must be read before their later reception.

\subsection{\emph{Antigone}: civic decree and a claimed higher duty}

In Sophocles' \emph{Antigone} 450--460, Antigone denies that Creon's proclamation can override ordinances that she attributes to the gods as unwritten and enduring (\sourceurl{https://www.perseus.tufts.edu/hopper/text?doc=Perseus\%3Atext\%3A1999.01.0186\%3Acard\%3D441}{Greek and English witness}). The scene dramatizes a conflict between civic office and a claimed higher obligation. It is character speech in a tragedy concerned with piety, kinship, rule, excess, and judgment, not a systematic theory of natural law.

The passage nevertheless attests that Greek drama could distinguish a civic decree from the justice of obeying it. Aristotle later cites Antigone's speech in his discussion of common law according to nature. The reception is explicit; the meaning Aristotle assigns to the example belongs to his own argument.

\subsection{Nature and convention require an account of justice}

The contrast between \emph{physis} and \emph{nomos} did not carry one moral conclusion. In Plato's \emph{Gorgias} 482e--484c, Callicles appeals to nature to defend domination by the stronger (\sourceurl{https://www.perseus.tufts.edu/hopper/text?doc=Perseus\%3Atext\%3A1999.01.0178\%3Atext\%3DGorg.\%3Apage\%3D482}{text}). He is an interlocutor whose position Socrates contests. The passage shows why an appeal to nature requires an account of what nature means and why it is normative.

Plato's \emph{Laws} begins by asking whether a constitution is attributed to a god or a human lawgiver (I, 624a--625a). Book X presents a materialist and conventionalist explanation of gods and justice before answering it through the priority of soul and intellect (888e--890a, 890d--892c, 896a--899d; \sourceurl{https://www.perseus.tufts.edu/hopper/text?doc=Perseus\%3Atext\%3A1999.01.0166\%3Abook\%3D10}{Book X}). The dialogue distinguishes divine attribution, rational cosmic order, and civic enactment. Its theology remains Greek philosophical and civic theology; any relation to Christian accounts must be demonstrated by later sources.

\subsection{Aristotle: natural justice, determination, and equity}

Aristotle gives the most explicit natural/legal distinction in this Greek corpus. In \emph{Nicomachean Ethics} V.7, 1134b18--1135a5, political justice is partly natural and partly legal. Legal justice includes matters that could initially be settled in more than one way but cease to be indifferent once authoritatively determined (\sourceurl{https://www.perseus.tufts.edu/hopper/text?doc=Perseus\%3Atext\%3A1999.01.0054\%3Abook\%3D5}{Book V}). Positive determination is therefore a constitutive part of political justice, not merely a departure from nature.

Aristotle also says that what is natural in political justice can vary in human affairs. At minimum, the text requires readers to distinguish the natural basis of justice from changing material circumstances and applications. It should not be assimilated without argument to a later scholastic account of immutable first principles.

\emph{Rhetoric} I.13 distinguishes a city's particular written or unwritten law from a common law according to nature and cites Antigone. The adjacent account of equity addresses cases in which legislation must speak generally but its wording does not fit an unforeseen case (\sourceurl{https://www.perseus.tufts.edu/hopper/text?doc=Perseus\%3Atext\%3A1999.01.0060\%3Abook\%3D1}{Book I, 1373b--1374b}). Equity seeks the reason of the rule; it is not a private power to suspend any inconvenient law.

In \emph{Politics} III.16, Aristotle contrasts rule by law with rule by appetite while recognizing that magistrates must decide what a general text cannot settle (\sourceurl{https://www.perseus.tufts.edu/hopper/text?doc=Perseus\%3Atext\%3A1999.01.0058\%3Abook\%3D3\%3Asection\%3D1287a}{1287a18--32}). Together, these passages preserve both governance by general law and prudent judgment in cases the text does not resolve.

\subsection{Stoicism: nature, reason, and common law}

Diogenes Laertius reports that the Stoic life according to nature conforms to a common law identified with right reason pervading all things and associated with Zeus's governance (VII.87--89; \sourceurl{https://www.perseus.tufts.edu/hopper/text?doc=Perseus\%3Atext\%3A1999.01.0258\%3Abook\%3D7\%3Achapter\%3D1}{text}). Cleanthes' \emph{Hymn to Zeus} expresses poetically a cosmos ordered through \emph{logos} and common law. A fragment of Chrysippus survives in Marcian's excerpt at Digest 1.3.2, where law is called ruler of divine and human affairs and a standard of justice (\sourceurl{https://droitromain.univ-grenoble-alpes.fr/Corpus/d-01.htm}{Digest Book 1}).

The evidence is fragmentary and mediated. Later Christian writers use related language, but direct reception must be shown source by source. Stoic immanent rational order is not identical with creation by a transcendent God, and conformity to nature is not the revealed New Law of grace.

\subsection{Result of the Greek survey}

\begin{threestable}{Question preserved}{Contribution of this corpus}{Historical limit}
Higher obligation & A civic command can be assessed under a claimed divine or common norm. & A dramatic claim or philosophical argument does not itself create a legal remedy.\\
Nature and convention & Nature can be invoked for rational order or for domination; the meaning and inference must be argued. & The terms do not carry a single theory across authors and genres.\\
Natural and legal justice & Political justice includes both nonconventional reasons and authoritative determinations. & Aristotle's account is not Aquinas's completed law taxonomy.\\
Equity and judgment & General rules require reasoned application where their wording does not fit the case. & Equity remains answerable to the rule's purpose and competent judgment.\\
Common rational law & Stoic sources connect nature, reason, divinity, and norm in a cosmopolitan framework. & Fragmentary transmission and different metaphysics limit claims of direct continuity.\\
\end{threestable}
